% NMI majority style: no "Introduction" heading (the opening runs unlabelled); the arXiv
% build keeps the heading.
\ifnmibuild\else
\section{Introduction}
\label{sec:intro}
\fi

Peer review is the gate through which scientific work passes, and it is under strain.
Submission volumes have grown faster than the reviewer pool, and the process is
demonstrably noisy: the NeurIPS consistency experiments found that a large fraction of
accept/reject decisions would flip under a re-run of the same process
\citep{cortes2021inconsistency,beygelzimer2023arbitrary}, with the asymmetry that
review is reliable at \emph{rejecting weak work} but poor at \emph{ranking strong
work}. Against this backdrop, LLM-based assistance is now being deployed at scale
\citep{thakkar2026feedback,biswas2026aaai}.

Most studies of automated review evaluate the \emph{text} a model produces: how well
its comments overlap a human reviewer's, how specific or factual they are
\citep{liang2024llm,darcy2024marg,yuan2022automate,du2024llmsassist}. For deployment
the more consequential question is whether the \emph{score} a system assigns is a valid
signal of the quality humans perceive. A score is what triages a queue or flags a
submission for a second look, so it must be validated as a measurement against an
external outcome.

% NMI majority style: prior work is folded into the opening section, here between the gap
% statement and the method. In the arXiv build this is skipped and Related Work is its own
% section (main.tex inputs 02_related_work separately).
\ifnmibuild\section{Related Work}
\label{sec:related}

\paragraph{Automated review, and grading by prompting.}
Using LLMs to assist peer review has moved from speculation to deployment: a randomized
study at ICLR~2025 improved review informativeness with an LLM feedback agent
\citep{thakkar2026feedback}, and AAAI~2026 ran an AI-review pilot at scale
\citep{biswas2026aaai}. The anchor empirical study found GPT-4 review comments overlap
human comments about as much as two humans overlap \citep{liang2024llm}; multi-agent and
pipeline systems raise specificity \citep{darcy2024marg,lu2024aiscientist}; recent
benchmarks measure LLM reviewer ability and its divergence from humans
\citep{li2026llmreviewer,reviewertoo2025}; and critical work documents brittleness,
hidden-prompt attacks, and inflated scores
\citep{du2024llmsassist,zhou2024reliable,lin2025hidden,zhou2025positivereview,russolatona2024lottery,pangram2026iclr},
prompting funding agencies and journals to restrict undisclosed AI review
\citep{science2025funding}. Almost all of this evaluates generated review \emph{text}.
Separately, \emph{grading by prompting} is an established paradigm: strong LLMs score
open-ended text against rubrics with near-human agreement (LLM-as-a-judge;
\citealp{zheng2023judging,liu2023geval}), and prompt-only scoring has been validated
against human grades in automated essay scoring \citep{mizumoto2023exploring}. We build
on that paradigm rather than claim it.

\paragraph{Noisy ground truth, and how to validate against it.}
Any claim that a signal tracks ``quality'' must be calibrated against the unreliability
of the human ground truth. The NeurIPS consistency experiments are canonical: roughly
half of accepted papers would change under a re-run, and review succeeds at rejecting
weak papers but struggles to rank strong ones
\citep{cortes2021inconsistency,beygelzimer2023arbitrary}; bias and its measurement are
well documented \citep{tomkins2017reviewerbias,stelmakh2019biases,shah2022challenges}.
For validating a score against noisy, ordinal outcomes the field uses rank correlation,
threshold-free AUROC, and calibration assessment \citep{spearman1904,guo2017calibration},
on corpora such as PeerRead and NLPeer \citep{kang2018peerread,dycke2023nlpeer};
\citet{checco2021aiassisted} validated a learned score against decisions and ratings,
though with a trained classifier rather than an LLM.

\paragraph{Positioning.}
The older paradigm for scholarly assessment \emph{trains} on reviews and decisions
(acceptance prediction; \citealp{kang2018peerread}) or \emph{fine-tunes} review
generators \citep{idahl2025openreviewer}. Prompt-only treatment of manuscripts has so
far targeted feedback \emph{text} \citep{liang2024llm} or has been validated only against
weak proxies, such as an author's own ratings of his own papers
\citep{thelwall2024chatgpt}. To our knowledge this is the first study to validate a
training-free, prompt-only manuscript quality \emph{score} against the public \emph{reject}
class and reviewer ratings, and to decompose how much of that validity comes from the
underlying model versus the engineered pipeline. The individual ingredients are not each
new --- prompting, pre-registration, and outcome validation all have precedent --- but
their combination on the observable reject class, with the model-versus-pipeline
reliability finding, is. The novelty is the validity result, not the prompting.
\fi

\textsc{AIPR} is an interactive platform for reviewing manuscripts: it reads a submitted
PDF and returns a structured, evidence-grounded review, of which a weighted overall
quality score across five dimensions is one component (\S\ref{sec:methods}). The score is
an AI first-pass signal, not a final human-reviewed verdict; AIPR is a human-in-the-loop
system in which a reviewer reads and can edit every part of the review, including the
passage-level highlights anchored on the manuscript, so the tool assists rather than
replaces the reviewer. AIPR grades by \emph{prompting alone}: it is not
fine-tuned on reviews, decisions, or ratings, so the score is a property of a prompted
frontier model, not of a classifier trained on outcomes. Here we isolate and benchmark
that first-pass score on first read, the hardest case for an automated signal and the one
a triage use relies on. Our ground truth is the public OpenReview record of a major
machine learning venue: each submission's decision tier (reject, poster, oral) and mean
reviewer rating; because the venue publishes \emph{rejected} submissions, the class our
central claim concerns is fully observable.

The paper turns on two pre-registered questions. \emph{Is the score valid?} That is, does
it agree with human outcomes well enough to flag weak submissions for a second look, the
low-end triage use we scope it to; we do not claim it predicts the acceptance decision,
ranks strong papers, or replaces a reviewer. \emph{What does the system add over the model
it runs on?} A frontier model is available to anyone, so we compare AIPR against a
one-paragraph prompt on the identical model and PDF, separating what the base model
contributes from what the engineered pipeline adds. The hypotheses and thresholds were
fixed before any score met any outcome (Appendix~\ref{app:prereg}); every number in the
paper is generated from the released data by one script, and the full study design is
diagrammed in Appendix~\ref{app:configs} (Fig.~\ref{fig:design}).
